\documentclass[10pt]{article}
\usepackage[margin=0.72in]{geometry}
\usepackage[T1]{fontenc}
\usepackage{lmodern}
\usepackage{microtype}
\usepackage{booktabs}
\usepackage{enumitem}
\usepackage{xcolor}
\usepackage{hyperref}
\usepackage{listings}
\usepackage{titlesec}
\hypersetup{colorlinks=true,urlcolor=blue,linkcolor=blue}
\titleformat{\section}{\large\bfseries}{\thesection}{0.5em}{}
\titleformat{\subsection}{\normalsize\bfseries}{\thesubsection}{0.5em}{}
\setlist{nosep,leftmargin=*}
\lstdefinestyle{code}{basicstyle=\ttfamily\scriptsize,breaklines=true,frame=single,numbers=left,numberstyle=\tiny\color{gray},showstringspaces=false,columns=fullflexible,tabsize=4}
\lstset{style=code}
\title{Pair-Programming Quick Reference\\\large German Credit Risk Case Study}
\author{Sizo Duma}
\date{}
\begin{document}
\maketitle
\begin{center}\small Practical TDD-first snippets for a 90-minute Python/SQL pairing session. This is a preparation aid, not a prediction of interview questions.\end{center}

\section{Interview approach: clarify, test, implement}
The case study asks for data preparation and exploration, three predictive features, credit-risk modelling, and validation. The submitted notebook is procedural and already uses a stratified 80/20 split (seed 42), a preprocessing pipeline, five-fold stratified CV, logistic regression and random forest. Keep the existing implementation as the baseline. Avoid rewriting the notebook or introducing classes unless the interviewers ask.

For each new requirement: clarify expected behavior and edge cases; write a failing test (Red); add the smallest implementation (Green); refactor while tests remain green; run targeted tests and summarize the change. Communicate trade-offs, especially around false positives, missed bad-risk applicants, probability calibration, and test-set leakage.

\section{Python: small testable functions}
These snippets assume the existing notebook variables \texttt{X}, \texttt{y}, \texttt{X\_train}, \texttt{X\_test}, \texttt{y\_train}, \texttt{y\_test}, \texttt{models}, \texttt{best}, \texttt{proba}, \texttt{pred}, \texttt{num\_cols}, and \texttt{cat\_cols}. Adapt imports and names to the checked-out notebook. Tests should be placed before implementation in a live TDD session.

\subsection{1. Cost threshold: make the convention explicit}
Here the positive label is \texttt{bad}. At threshold $t$, an applicant is flagged when $p(\mathrm{bad})\geq t$. This illustrative loss assigns cost to approving a bad applicant and rejecting a good applicant. The numbers are assumptions for a coding exercise, not estimated bank economics.

\begin{lstlisting}[language=Python,caption={Test first: minimum expected classification cost}]
import numpy as np

def test_cost_threshold_selects_lower_cost_candidate():
    p = np.array([0.1, 0.4, 0.8])
    y = np.array([0, 1, 1])
    result = choose_cost_threshold(
        y, p, thresholds=[0.2, 0.6],
        cost_approve_bad=5, cost_reject_good=1,
    )
    assert result["threshold"] in {0.2, 0.6}
    assert result["cost"] == min(
        expected_cost(y, p >= 0.2, 5, 1),
        expected_cost(y, p >= 0.6, 5, 1),
    )
\end{lstlisting}

\begin{lstlisting}[language=Python,caption={Implementation: evaluate decisions and select among candidate thresholds}]
def expected_cost(y_true, flagged_bad, cost_approve_bad, cost_reject_good):
    y_true = np.asarray(y_true, dtype=int)
    flagged_bad = np.asarray(flagged_bad, dtype=bool)
    false_negative = (y_true == 1) & ~flagged_bad
    false_positive = (y_true == 0) & flagged_bad
    return float(cost_approve_bad * false_negative.sum()
                 + cost_reject_good * false_positive.sum())

def choose_cost_threshold(y_true, probability, thresholds,
                          cost_approve_bad, cost_reject_good):
    if cost_approve_bad < 0 or cost_reject_good < 0:
        raise ValueError("Costs must be non-negative.")
    candidates = []
    for threshold in thresholds:
        decisions = np.asarray(probability) >= threshold
        candidates.append({
            "threshold": float(threshold),
            "cost": expected_cost(y_true, decisions,
                                  cost_approve_bad, cost_reject_good),
        })
    if not candidates:
        raise ValueError("Provide at least one threshold.")
    return min(candidates, key=lambda row: row["cost"])
\end{lstlisting}

Choose the threshold using training-only validation predictions, then evaluate it once on the untouched holdout. A 5:1 cost ratio is illustrative only. Real expected loss also involves exposure, loss given default, recoveries, pricing, policy constraints, and customer impact. Do not equate this simple loss with a bank's complete economic objective.

\subsection{2. Feature engineering without leakage}
Feature engineering should happen inside the preprocessing pipeline so every CV fold learns transformations from its training fold. A deterministic row-wise transformation can use a \texttt{FunctionTransformer}; test the transformation independently first.

\begin{lstlisting}[language=Python,caption={Test first: installment burden is amount divided by duration}]
import pandas as pd

def test_add_risk_features_calculates_burden_and_age_band():
    sample = pd.DataFrame({"credit_amount": [1200], "duration": [12],
                           "age": [24]})
    result = add_risk_features(sample)
    assert result.loc[0, "monthly_installment_burden"] == 100
    assert result.loc[0, "age_band"] == "under_25"
\end{lstlisting}

\begin{lstlisting}[language=Python,caption={Implementation: deterministic feature derivation}]
def add_risk_features(frame):
    result = frame.copy()
    result["monthly_installment_burden"] = (
        result["credit_amount"] / result["duration"].replace(0, np.nan)
    )
    result["age_band"] = pd.cut(
        result["age"], bins=[0, 25, 40, np.inf],
        labels=["under_25", "25_to_40", "over_40"], right=False,
    ).astype("string")
    return result
\end{lstlisting}

Before adding these fields, consider whether they are economically meaningful and whether they duplicate existing predictors. Age-derived features also warrant fairness and legal review. Do not present predictive associations as causal drivers.

\subsection{3. Pipeline factory and testable preprocessing}
A small function can create the existing pipeline without adding an abstraction layer. Any new transformer must be fitted within the pipeline and CV fold.

\begin{lstlisting}[language=Python,caption={Test first: pipeline exposes probabilities and preserves preprocessing}]
def test_make_model_pipeline_predicts_probabilities():
    pipeline = make_model_pipeline(
        LogisticRegression(max_iter=2000), num_cols, cat_cols
    )
    pipeline.fit(X_train, y_train)
    probability = pipeline.predict_proba(X_test.iloc[:3])[:, 1]
    assert probability.shape == (3,)
    assert np.isfinite(probability).all()
    assert ((probability >= 0) & (probability <= 1)).all()
\end{lstlisting}

\begin{lstlisting}[language=Python,caption={Implementation: pipeline factory}]
def make_model_pipeline(estimator, numeric_columns, categorical_columns):
    numeric = Pipeline([
        ("imputer", SimpleImputer(strategy="median")),
        ("scale", StandardScaler()),
    ])
    categorical = Pipeline([
        ("imputer", SimpleImputer(strategy="most_frequent")),
        ("onehot", OneHotEncoder(handle_unknown="ignore")),
    ])
    preprocess = ColumnTransformer([
        ("numeric", numeric, numeric_columns),
        ("categorical", categorical, categorical_columns),
    ])
    return Pipeline([("prep", preprocess), ("model", estimator)])
\end{lstlisting}

\subsection{4. Evaluation helper}
Keep ranking metrics on probabilities and threshold metrics on labels. Brier score is a probability error measure; it is not the same as ROC-AUC.

\begin{lstlisting}[language=Python,caption={Test first: evaluation returns expected metrics}]
def test_evaluate_binary_model_reports_auc_and_brier():
    y_true = np.array([0, 0, 1, 1])
    p = np.array([0.1, 0.4, 0.6, 0.9])
    result = evaluate_binary_model(y_true, p, threshold=0.5)
    assert 0 <= result["roc_auc"] <= 1
    assert result["brier"] >= 0
    assert result["threshold"] == 0.5
\end{lstlisting}

\begin{lstlisting}[language=Python,caption={Implementation: compact evaluation summary}]
def evaluate_binary_model(y_true, probability, threshold=0.5):
    prediction = (np.asarray(probability) >= threshold).astype(int)
    return {
        "roc_auc": roc_auc_score(y_true, probability),
        "average_precision": average_precision_score(y_true, probability),
        "balanced_accuracy": balanced_accuracy_score(y_true, prediction),
        "brier": brier_score_loss(y_true, probability),
        "threshold": float(threshold),
        "confusion_matrix": confusion_matrix(y_true, prediction),
    }
\end{lstlisting}

\subsection{5. Input validation and single-applicant scoring}
\texttt{OneHotEncoder(handle\_unknown="ignore")} avoids an encoder exception for unseen categories; it does not mean an unfamiliar category is valid. Validate inputs against a documented schema and defensible business constraints.

\begin{lstlisting}[language=Python,caption={Test first: reject a missing feature}]
import pytest

def test_validate_applicant_rejects_missing_columns():
    with pytest.raises(ValueError, match="Missing columns"):
        validate_applicants(pd.DataFrame({"age": [30]}), X.columns)
\end{lstlisting}

\begin{lstlisting}[language=Python,caption={Implementation: schema guard and probability scoring}]
def validate_applicants(rows, expected_columns):
    missing = set(expected_columns) - set(rows.columns)
    extra = set(rows.columns) - set(expected_columns)
    if missing:
        raise ValueError(f"Missing columns: {sorted(missing)}")
    if extra:
        raise ValueError(f"Unexpected columns: {sorted(extra)}")
    if rows[list(expected_columns)].isna().any().any():
        raise ValueError("Applicant data contains missing values.")
    return rows.loc[:, expected_columns].copy()

def score_applicants(rows, fitted_pipeline, expected_columns):
    checked = validate_applicants(rows, expected_columns)
    return fitted_pipeline.predict_proba(checked)[:, 1]
\end{lstlisting}

Do not return fabricated ``top risk drivers'' from a scoring function. To explain a prediction, use a defined, tested explanation method and label its limitations. Validate numeric ranges only when confirmed by the product/data contract; do not infer them from the observed sample minimum and maximum.

\section{SQL quick reference for analytics}
The referenced Packt repository, \emph{SQL for Data Analytics, Fourth Edition}, covers data loading, filtering, aggregation, CTEs, window functions, performance, and inferential statistics. Examples below use mostly standard SQL, with PostgreSQL-specific syntax labeled. Confirm the target database dialect before using features such as \texttt{FILTER}, casts, or \texttt{COPY}.

Assume a table named \texttt{credit\_applicants} and a quoted target column \texttt{"class"}. In production, use parameterized queries from Python rather than concatenating untrusted values into SQL.

\subsection{Profile labels and bad rate}
\begin{lstlisting}[language=SQL,caption={ANSI-style summary with conditional aggregation}]
SELECT
    COUNT(*) AS applicants,
    SUM(CASE WHEN "class" = 'good' THEN 1 ELSE 0 END) AS good_count,
    SUM(CASE WHEN "class" = 'bad' THEN 1 ELSE 0 END) AS bad_count,
    AVG(CASE WHEN "class" = 'bad' THEN 1.0 ELSE 0.0 END) AS bad_rate
FROM credit_applicants;
\end{lstlisting}

\subsection{Bad rates by category}
\begin{lstlisting}[language=SQL,caption={Group counts and observed bad rate}]
SELECT
    checking_status,
    COUNT(*) AS applicants,
    SUM(CASE WHEN "class" = 'bad' THEN 1 ELSE 0 END) AS bad_count,
    AVG(CASE WHEN "class" = 'bad' THEN 1.0 ELSE 0.0 END) AS bad_rate
FROM credit_applicants
GROUP BY checking_status
ORDER BY bad_rate DESC;
\end{lstlisting}

Rates should be presented with group counts; a small category may have a volatile rate. This descriptive table is not a causal analysis.

\subsection{Basic data-quality checks}
\begin{lstlisting}[language=SQL,caption={Check nulls and obvious numeric violations}]
SELECT
    COUNT(*) AS row_count,
    SUM(CASE WHEN "class" IS NULL THEN 1 ELSE 0 END) AS missing_target,
    SUM(CASE WHEN duration IS NULL THEN 1 ELSE 0 END) AS missing_duration,
    SUM(CASE WHEN credit_amount IS NULL THEN 1 ELSE 0 END) AS missing_amount,
    SUM(CASE WHEN duration <= 0 THEN 1 ELSE 0 END) AS nonpositive_duration,
    SUM(CASE WHEN credit_amount <= 0 THEN 1 ELSE 0 END) AS nonpositive_amount
FROM credit_applicants;
\end{lstlisting}

\subsection{Age-group descriptive slice}
\begin{lstlisting}[language=SQL,caption={Observed label rates by age bracket (PostgreSQL CASE is portable)}]
SELECT
    CASE WHEN age < 25 THEN 'under_25'
         WHEN age > 40 THEN 'over_40'
         ELSE '25_to_40' END AS age_group,
    COUNT(*) AS applicants,
    SUM(CASE WHEN "class" = 'bad' THEN 1 ELSE 0 END) AS bad_count,
    AVG(CASE WHEN "class" = 'bad' THEN 1.0 ELSE 0.0 END) AS bad_rate
FROM credit_applicants
GROUP BY 1
ORDER BY 1;
\end{lstlisting}

This describes historical label rates, not model false-positive rates. For false-positive rates, query actual and predicted labels from a holdout predictions table, and report only as a limited diagnostic.

\subsection{Window-function example}
\begin{lstlisting}[language=SQL,caption={Rank categories by observed bad rate}]
WITH category_rates AS (
    SELECT checking_status,
           COUNT(*) AS applicants,
           AVG(CASE WHEN "class" = 'bad' THEN 1.0 ELSE 0.0 END) AS bad_rate
    FROM credit_applicants
    GROUP BY checking_status
)
SELECT checking_status, applicants, bad_rate,
       DENSE_RANK() OVER (ORDER BY bad_rate DESC) AS risk_rate_rank
FROM category_rates
ORDER BY risk_rate_rank, checking_status;
\end{lstlisting}

Use SQL for data access and descriptive summaries; preserve the notebook's Python pipeline for leakage-safe model fitting and evaluation. Never select model features or thresholds from holdout-derived SQL summaries.

\section{Talking points: explain without overselling}
\begin{itemize}
  \item The target is binary; treating \texttt{bad} as positive makes default detection metrics easier to interpret.
  \item Stratification preserves approximate label proportions across train/test partitions and CV folds; it does not remove uncertainty from a small dataset.
  \item ROC-AUC measures ranking, not whether predicted probabilities are accurate or whether a lending policy is profitable.
  \item Precision and recall expose different mistakes: missed bad applicants versus good applicants flagged for review or rejected.
  \item A class-weighted model can change probability interpretation; compare calibration and base rates, not AUC alone.
  \item Permutation importance measures the change in predictive score after shuffling a feature. It is not causal and can be affected by correlated predictors.
  \item This historical dataset is a teaching example, not a production underwriting dataset. Real use would require governance, calibration, fairness review, monitoring, policy controls, and legal review.
\end{itemize}

\section{Reference}
Jun Shan, Haibin Li, Matt Goldwasser, Upom Malik, and Benjamin Johnston, \emph{SQL for Data Analytics, Fourth Edition}, Packt Publishing. Companion code repository: \url{https://github.com/PacktPublishing/SQL-for-Data-Analytics-Fourth-Edition}. The book repository's README lists examples covering table creation, COPY, Python data manipulation, SELECT, transformations, aggregation, window functions, performance, JSON/arrays, inferential statistics, and an analytics case study.

\end{document}
